Generative Incompletion is what happens when something genuine escapes 
before it has fully resolved --- when the internal guard has lowered enough, 
or the feeling has grown present enough, that expression arrives ahead 
of articulation.

It is not a technique. It cannot be aimed for. The moment a man decides 
to be incomplete, he has already closed himself --- the decision is its own 
form of construction. What this section describes is not a way of 
manufacturing openness, but a discipline of \textit{not closing what is 
already open}.

This is why it is the most potent emergence of genuine presence: it is 
the point at which presence becomes \textit{visible} --- not because the 
man chose to reveal something, but because something real became too 
present to contain.

When a genuine fragment lands --- meaning begun but not resolved --- 
the receiver's mind moves automatically to complete it. The meaning they 
generate comes from inside them. They does not receive it. They build it. 
This is the mechanism behind what will later be examined as 
\hyperref[sec:co_authorship]{Co-Authorship}: investment produced not 
by what was delivered, but by what they was invited to finish.